\section{Uniform Laplace envelopes and odd-population extremizers}\label{sec:odd_uniform}
Write $g(d)=d/\operatorname{artanh}d$ with $g(0)=1$, and recall $L_N=[(N-1)\operatorname{artanh}(1/N)]^{-1}$. The Laplace comparison below applies to both parities; only the final optimization over discrete success counts uses oddness.

\begin{theorem}[Parity-uniform Bernoulli envelope]\label{odd:thm:envelope}
Let $N\ge7$ be an integer with $3\le m\le k\le\lfloor N/2\rfloor$ and $N>2m$, and let $H\sim\Hyp(N,k,m)$. Write $\psi=\psi_{N,k,m}$ and set
\[
 p=\frac{k}{N},\quad q=1-p,\quad d=1-2p,\quad
 M=\frac{m(N-m)}{N-1},\quad a=\frac{N-2m}{N-2}>0.
\]
For every $t\in\R$,
\begin{equation}\label{odd:eq:laplace}
 \boxed{\quad
 \psi(t)\le\frac{M}{a^2}\bigl[\log(q+pe^{at})-pat\bigr].\quad}
\end{equation}
The Laplace inequality is strict for $t\ne0$. If $k<N/2$, then
\begin{equation}\label{odd:eq:BernEnvelope}
 \sigma^2_{\rm opt}(H)<M\sigma^2_{\rm opt}(\Bern(p)),\qquad
 \frac{4N}{m(N-m)}\sigma^2_{\rm opt}(H)
 <\frac{N}{N-1}g(d).
\end{equation}
If $N$ is even and $k=N/2$, the two optimal proxies are equal: both sides of the first comparison equal $M/4$, and both sides of the normalized comparison equal $N/(N-1)$. For odd $N$, $d\ge 1/N$, so the last expression in \eqref{odd:eq:BernEnvelope} is at most $L_N$.
All these estimates are finite-population statements.
\end{theorem}

\subsection{Tilted mean and variance as a curve}
Throughout the envelope proof assume the hypotheses of Theorem~\ref{odd:thm:envelope}.
Let $\E_t$ and $\Var_t$ denote exponential tilting by $e^{tH}$, and write
\[
 u(t)=\E_tH,\qquad v(t)=\Var_t(H).
\]
Every mass on $\{0,\ldots,m\}$ is positive. Hence $u'(t)=v(t)>0$ and $u$ maps $\R$ bijectively onto $(0,m)$. Thus $v$ defines a smooth function $V(u)$. At the endpoints $V(u)\to0$, and, writing $\mu=mp$,
\begin{equation}\label{odd:eq:variance0}
 V(\mu)=Mpq=\frac M4(1-d^2).
\end{equation}

\begin{lemma}[Real-root admissibility]\label{odd:lem:PB}
For every $0<u<m$,
\begin{equation}\label{odd:eq:PB}
 0<V(u)\le\frac{u(m-u)}m.
\end{equation}
\end{lemma}
\begin{proof}
Apply the Jacobi factorization and Bernoulli representation from Section \ref{sec:jacobi} with degree $m$ and parameters $N-k-m$ and $k-m$. If the tilted Bernoulli probabilities are $p_1,\ldots,p_m$, then
\[
 u=\sum_i p_i,\qquad V=\sum_i p_i(1-p_i)
 =u-\sum_i p_i^2\le u-\frac{u^2}{m}.
\]
\end{proof}

Define
\[
 A(u)=u(u+N-k-m),\qquad C(u)=(k-u)(m-u)
\]
and the polynomial
\begin{equation}\label{odd:eq:F}
 F(u,v)=-(N-1)v^2+[A+C-A'C+C'A]v+AC.
\end{equation}
Primes on $A,C$ denote differentiation with respect to $u$.

\begin{lemma}[Differential identity]\label{odd:lem:ode}
Writing $y=u-\mu$, the true variance curve obeys
\begin{equation}\label{odd:eq:curveODE}
 F(u,V(u))+NyV(u)V'(u)=0.
\end{equation}
\end{lemma}
\begin{proof}
Let $z=e^t$. The tilted adjacent-mass relation, summed over the support, gives
\[
 \E_t[H(H+N-k-m)]=z\E_t[(k-H)(m-H)].
\]
Consequently
\[
 z=\frac{A(u)+V(u)}{C(u)+V(u)}.
\]
Both numerator and denominator are positive. Differentiate its logarithm with respect to $u$, using $dt/du=1/V$:
\[
 \frac{A'+V'}{A+V}-\frac{C'+V'}{C+V}=\frac1V.
\]
Multiplication and collection yield $(C-A)VV'=F(u,V)$. Since $C-A=km-Nu=-Ny$, this is \eqref{odd:eq:curveODE}. The identity also holds at $y=0$.
\end{proof}

\subsection{A quadratic envelope and an all-parameter certificate}
Define a quadratic polynomial in $u$ by
\begin{equation}\label{odd:eq:P}
 P(u)=\frac M4(1-d^2)+ad(u-\mu)-\frac{a^2}{M}(u-\mu)^2.
\end{equation}
The envelope depends quadratically on the tilted mean.

\begin{lemma}[Quadratic variance envelope]\label{odd:lem:barrier}
For every $0<u<m$,
\begin{equation}\label{odd:eq:barrier}
 \boxed{\quad V(u)\le P(u),\quad}
\end{equation}
with equality if and only if $u=\mu$.
\end{lemma}

For the variance comparison, put
\[
 z_0=\frac{2(u-\mu)}M,\quad
 G=\frac{N^2(m-1)(N-m-1)}{4(N-2)^2(N-1)}>0
\]
and, for free variables $d,z_0$, define
\begin{equation}\label{odd:eq:W}
 W(d,z_0)=N-4+3Nd^2-4Nad z_0+[(N+1)a^2-M]z_0^2.
\end{equation}
An exact substitution into \eqref{odd:eq:F} gives
\begin{equation}\label{odd:eq:residual}
 \boxed{\quad F(u,P(u))+N(u-\mu)P(u)P'(u)
 =-G(u-\mu)^2W(d,z_0).\quad}
\end{equation}
At a noncentral contact $V=P$, the sign of the residual is determined by $W$.

Set
\begin{align*}
 \delta&=\frac{N-2m}{N},&
 b&=\frac{N-m}{N-2},&
 c&=\frac{(N-m)[3N^2-4N-4m(N-1)]}{(N-1)(N-2)^2},\\
 E(d,z_0)&=1-d^2+4bdz_0-cz_0^2.
\end{align*}
The exact admissibility identity is
\begin{equation}\label{odd:eq:E}
 \frac{u(m-u)}m-P(u)=\frac{m(m-1)}{4(N-1)}E(d,z_0).
\end{equation}
Thus, at any contact $V=P$, Lemma~\ref{odd:lem:PB} implies $E\ge0$.
Also $m\le k\le N/2$ implies $0\le d\le\delta$; equality $d=0$ occurs precisely at the even balanced margin $k=N/2$.

Define
\begin{equation}\label{odd:eq:lambda}
 \lambda=\frac{m^2-mN+3m+N^2-4N}{2(N-m)},\qquad
 \alpha=\frac{N-4-\lambda}{\delta^2},
\end{equation}
and
\[
 A_0=3N+\lambda+\alpha,\qquad
 B_0=-2Na-2\lambda b,\qquad
 C_0=(N+1)a^2-M+\lambda c.
\]
The following is a polynomial identity after clearing the positive denominators:
\begin{equation}\label{odd:eq:certificate}
 \boxed{\quad
 W=\lambda E+\alpha(\delta^2-d^2)
   +A_0d^2+2B_0dz_0+C_0z_0^2.\quad}
\end{equation}
Moreover,
\begin{equation}\label{odd:eq:det}
 A_0C_0-B_0^2
 =\frac{N^2U_0V_0}{(N-2m)^2(N-2)^2(N-1)},
\end{equation}
where
\begin{align*}
 U_0&=N^2-2N-m^2-mN+3m,\\
 V_0&=m^3-6m^2-2mN^2+14mN-3m+N^3-7N^2+6N.
\end{align*}

For a direct all-degree sign verification write
\[
 j=m-3\ge0,\qquad r=N-2m\ge1.
\]
Then
\begin{align}
 2(N-m)\lambda
 &=3j^2+3jr+13j+r^2+5r+12>0,\label{odd:eq:lambda_pos}\\
 2(N-m)(N-4-\lambda)
 &=j^2+3jr+3j+r^2+5r>0,\label{odd:eq:alpha_pos}\\
 U_0&=j^2+3jr+5j+r^2+7r+6>0,\label{odd:eq:U_pos}\\
 V_0&=j^3+4j^2r+3j^2+4jr^2+10jr+r^3+5r^2>0.\label{odd:eq:V_pos}
\end{align}
It follows that $\lambda,\alpha>0$, $A_0>0$, and $A_0C_0-B_0^2>0$.
The quadratic form in \eqref{odd:eq:certificate} is positive definite. Therefore
\begin{equation}\label{odd:eq:Wpositive}
 E\ge0,\quad 0\le d\le\delta\quad\Longrightarrow\quad W(d,z_0)>0.
\end{equation}
The coefficient identities establish strict positivity for every allowed $N,m$.

\begin{proof}[Proof of Lemma~\ref{odd:lem:barrier}]
At the two endpoints the envelope is positive. Explicitly,
\begin{align}
 P(0)&=\frac{km(m-1)[(N-m)(N-2)-k(m-1)]}
 {(N-m)(N-2)^2(N-1)}>0,\label{odd:eq:end0}\\
 P(m)&=\frac{m(N-k)(m-1)[k(m-1)+(N-2m)(N-1)]}
 {(N-m)(N-2)^2(N-1)}>0.\label{odd:eq:endm}
\end{align}
For the first sign, $k\le N-m$ bounds the bracket below by
$(N-m)(N-m-1)>0$. The second sign is immediate.
Let $f(u)=V(u)-P(u)$, extending continuously to $[0,m]$.
Then $f(0)<0$, $f(m)<0$, and \eqref{odd:eq:variance0} gives $f(\mu)=0$.

At any zero $u\ne\mu$ of $f$, we have $P=V>0$. Combining
\eqref{odd:eq:curveODE} with \eqref{odd:eq:residual} gives
\begin{equation}\label{odd:eq:crossing}
 f'(u)=\frac{G(u-\mu)W(d,z_0)}{NP(u)}.
\end{equation}
At this zero, \eqref{odd:eq:E} and \eqref{odd:eq:PB} imply $E\ge0$, so
\eqref{odd:eq:Wpositive} applies. Thus every zero left of $\mu$ has strictly negative derivative, and every zero right of $\mu$ has strictly positive derivative.

If $f$ were positive somewhere in $(0,\mu)$, the left endpoint of a component of $\{f>0\}$ would lie in $(0,\mu)$, and would be a zero with nonnegative derivative, a contradiction. If $f$ were positive somewhere in $(\mu,m)$, the right endpoint of a positive component would give a zero with nonpositive derivative, again a contradiction. Hence $f\le0$ everywhere. Any additional interior zero would be a local maximum of $f$ and have derivative zero, contradicting \eqref{odd:eq:crossing}. This proves \eqref{odd:eq:barrier} and its equality statement.
\end{proof}

\subsection{Integration gives the scaled Bernoulli Laplace bound}
Write $y(t)=u(t)-\mu$. Lemma~\ref{odd:lem:barrier} says
\[
 y'(t)\le \mathcal P(y(t)),\qquad
 \mathcal P(y)=Mpq+ady-\frac{a^2}{M}y^2,
\]
with strict inequality for $t\ne0$. The exact solution starting at zero to
$Y'=\mathcal P(Y)$ is
\begin{equation}\label{odd:eq:Y}
 Y(t)=\frac{M}{a}\left[\frac{pe^{at}}{q+pe^{at}}-p\right].
\end{equation}
For the two-sided comparison, put
\[
 \Delta(t)=\mathcal P(y(t))-y'(t),\qquad
 h(t)=ad-\frac{a^2}{M}(y(t)+Y(t)).
\]
Then $\Delta(t)>0$ for $t\ne0$, and $e(t)=y(t)-Y(t)$ solves
\[
 e'=he-\Delta,\quad e(0)=0,\qquad
 e(t)=-e^{\int_0^t h}\int_0^t e^{-\int_0^s h}\Delta(s)\,ds.
\]
Thus $y(t)<Y(t)$ for $t>0$ and $y(t)>Y(t)$ for $t<0$.
Since $\psi'(t)=y(t)$ and $\psi(0)=0$, integration, with the correct orientation for negative $t$, yields
\[
 \psi(t)<\int_0^tY(s)\,ds
 =\frac{M}{a^2}[\log(q+pe^{at})-pat]\quad(t\ne0).
\]
This proves \eqref{odd:eq:laplace}.

The exact Bernoulli proxy is
\begin{equation}\label{odd:eq:bern}
 \sigma^2_{\rm opt}(\Bern(p))=\frac{q-p}{2\log(q/p)}
 =\frac{d}{4\artanh d}.
\end{equation}
On substituting $s=at$ in the supremum, \eqref{odd:eq:laplace} implies
$\sigma^2_{\rm opt}(H)\le M\sigma^2_{\rm opt}(\Bern(p))$.
To optimize the comparison, extend the hypergeometric quotient continuously to zero. It tends to zero at either infinity, so its maximum is attained at a finite point. At every nonzero tilt it is strictly below the Bernoulli upper envelope. If $d>0$, its value at zero is $Mpq$, strictly below $M\sigma^2_{\rm opt}(\Bern(p))=Md/(4\artanh d)$, because
\[
 \artanh d=\int_0^d\frac{dx}{1-x^2}<\frac{d}{1-d^2}\qquad(0<d<1).
\]
Thus the proxy comparison is strict for $k<N/2$. If $d=0$, the integrated Laplace bound and $\log\cosh x\le x^2/2$ give $\sigma^2_{\rm opt}(H)\le M/4$. The variance at zero tilt is $Mpq=M/4$, so equality holds in the optimal-proxy comparison. Finally, $g(d)$ strictly decreases for $d>0$ by the same integral inequality. In the odd case $d\ge1/N$, and hence $g(d)\le g(1/N)$. This completes Theorem~\ref{odd:thm:envelope} for both parities.

\subsection{Exact odd-population extremizers}
\begin{theorem}[Odd-population hypergeometric extremizers]\label{odd:thm:odd}
For every odd $N\ge3$,
\begin{equation}\label{odd:eq:exact}
 \boxed{\quad
 \max_{1\le k,m<N}\frac{4N}{m(N-m)}\sigma^2_{\mathrm{opt}}(H_{N,k,m})
 =\frac{1}{(N-1)\operatorname{artanh}(1/N)}.\quad}
\end{equation}
Equality holds precisely for $m\in\{1,N-1\}$ and $k\in\{(N-1)/2,(N+1)/2\}$.
\end{theorem}
\begin{proof}
Complementation of either margin replaces the centered variable by its negative. It leaves the optimal proxy unchanged, and sample complementation preserves $m(N-m)$. Reduce first to $1\le k,m\le(N-1)/2$.

For $m=1$, the Bernoulli formula \eqref{eq:bernproxy} gives
\[
 \frac{4N}{N-1}\sigma^2_{\mathrm{opt}}(H_{N,k,1})
 =\frac{N}{N-1}g(1-2k/N)\le L_N,
\]
with equality exactly at $k=(N-1)/2$. For $2\le m\le k$, Corollary \ref{cor:two_bern_envelope} and Theorem \ref{odd:thm:envelope} give strict inequality.

If $k<m$, transposition gives
\[
 \frac{4N}{m(N-m)}\sigma^2_{\mathrm{opt}}(H_{N,k,m})
 =\frac{k(N-k)}{m(N-m)}
 \frac{4N}{k(N-k)}\sigma^2_{\mathrm{opt}}(H_{N,m,k})<L_N,
\]
because the first factor is strictly below one and the transposed law is in the canonical triangle. Restoring the complements gives the stated equality cases. For $N=3$ only the singleton layer occurs; for $N=5$ the singleton and two-draw layers exhaust the canonical triangle.
\end{proof}

Applying \eqref{eq:LKidentity} proves the odd formula in Theorem \ref{thm:exchange}. The comparison is finite-population and covers every sample size, including sequences with $m_N\to\infty$, whether $m_N=o(N)$ or $m_N/N$ stays bounded away from zero.
